\begin{abstract}
  \label{sec:abstract}

  As machine learning expands into diverse domains, TinyML has emerged as a crucial paradigm for deploying models on highly resource-constrained microcontrollers, which typically feature less than 256~KB of RAM. However, executing complex mathematical operations on these devices remains a significant challenge, necessitating novel model designs and hardware-aware optimization.
  The Mamba architecture, built around State-Space Model, is a promising candidate due to its compact parameterization and strong performance on long-context tasks. Nevertheless, Mamba was originally designed for highly parallelized GPUs, making its adaptation for TinyML non-trivial. This paper evaluates Mamba deployment strategies on microcontrollers using TensorFlow Lite Micro.
  We propose architecture modifications and optimization techniques tailored specifically to microcontroller constraints. Our deployment of a quantized Mamba model achieves a 60.4~KB peak RAM footprint on a Keyword Spotting task, a 74\% memory reduction compared to state-of-the-art work (MambaLite-Micro). Furthermore, we analyze the trade-offs of quantization, demonstrating that while it substantially reduces memory, it can introduce latency overhead on hardware lacking acceleration of INT8 operations.
  To mitigate code size and loop-unrolling overheads, we introduce a model-splitting technique that enables the execution of larger models. Our findings demonstrate that while Mamba is a viable architecture for TinyML, further research is required to fully optimize State Space Model implementations for edge hardware.

\end{abstract}